\subsection{A larger pendant corner attached through a projective endpoint}

\begin{proposition}[Proved]
\label{CP2:A-pendant-eleven-permitted-triangle}
Let $C$ have four objects and acyclic support, and suppose
an object $p$ is strictly adjacent only to $q$.
Assume
\[
 |\operatorname{End}_C(q)|\leq2,
 \qquad |\operatorname{End}_C(p)|\leq11.
\]
Suppose the retained full subcategory $E=C\setminus\{p\}$
is a permitted triangle, in the sense defined before
Theorem~\ref{CP2:A-pendant-edge-permitted-triangle}.
If $BC$ has bounded rational homology and
$\dim H^2(BC;\mathbb Q)\leq1$, then every rational
degree-two cohomology product on $BC$ vanishes.
There is no restriction on the number of arrows between
$p$ and $q$.
\end{proposition}
\begin{proof}
Put $F=C|_{\{p,q\}}$ and let $I$ be the corner category
at $q$. Every simplex of $NC$ lies in $NE$ or $NF$.
Indeed, a simplex containing $p$ and a third distinct
object would contain a strict chain whose composite
gives an arrow between $p$ and that third object,
contrary to the pendant condition. Thus the nerve is
the actual union $NE\cup_{NI}NF$.

Every monoid of order at most two has semisimple
rational algebra. In particular $BI$ is rationally
acyclic and its trivial rational modules are projective
on both sides. Lemma~\ref{CP2:A-low-homology-gluing}
therefore passes bounded rational homology and the
second Betti bound to both pieces. The permitted
triangle statements preceding
Theorem~\ref{CP2:A-pendant-edge-permitted-triangle}
annihilate all degree-two products on $BE$.

Within the two-object category $F$, the endpoint $q$
is a source or sink. Apply the actual orbit functor of
Lemma~\ref{CP2:A-high-orbits} at that endpoint.
The appropriate trivial module is projective, so this
functor gives an isomorphism of rational cohomology
rings and a rational homology equivalence from $BF$
to the ordinary act cone over
$\operatorname{End}_C(p)$, or its opposite.
The orbit set of the nonempty stalk is nonempty.
The cone inherits bounded rational homology from $BF$;
Theorem~\ref{CP2:A-bounded-eleven-acts} consequently
makes all its degree-two products zero. Pullback along
the specified functor does the same on $BF$.

Finally the same gluing lemma makes restriction in
degree four injective. A product that restricts to zero
on both $BE$ and $BF$ is therefore zero on $BC$.
The argument uses an actual nerve union and an actual
orbit quotient; it does not replace the pendant by a
homotopy-equivalent monoid inside the original category.
\end{proof}

\begin{remark*}
The frozen four-object eighteen-arrow rows $1491$ and
$1492$ have $p=0$, $q=3$, and pendant corner order seven.
Their retained triangles have profiles respectively
$(1,1,1;2,2,1)$ and $(1,1,2;2,1,1)$, so both satisfy
the proposition. The two attachment corners have orders
one and two, respectively.
\end{remark*}
